\section*{Chapter \ref{ch:g1:light-and-shadows} --- Light and Shadows}

\begin{solution}{exo:g1:light-and-shadows:1}
Sources: the candle flame, the firefly, the torch. Non-sources: the
Moon, the mirror, the white wall --- they only send back \omterm{def:g1:light-and-shadows:source}{light} that
falls on them.
\end{solution}

\begin{solution}{exo:g1:light-and-shadows:2}
\omterm{def:g1:light-and-shadows:source}{Light} is missing. There is no source in the room, so no \omterm{def:g1:light-and-shadows:source}{light} falls
on the things, and things without \omterm{def:g1:light-and-shadows:source}{light} on them cannot be seen.
\end{solution}

\begin{solution}{exo:g1:light-and-shadows:3}
A \omterm{def:g1:light-and-shadows:source}{light source}, a blocker, and a wall (or screen, or ground) --- in
that order along a straight line: the source shines, the blocker
blocks, the wall shows the dark patch.
\end{solution}

\begin{solution}{exo:g1:light-and-shadows:4}
The \omterm{def:g1:light-and-shadows:shadow}{shadow} is on your right --- opposite the Sun. After turning right
around, the Sun is on your right, so the \omterm{def:g1:light-and-shadows:shadow}{shadow} lies on your left:
always on the far side from the Sun.
\end{solution}

\begin{solution}{exo:g1:light-and-shadows:5}
Moving the toy \emph{toward the lamp} makes the \omterm{def:g1:light-and-shadows:shadow}{shadow} huge; moving
it \emph{toward the wall} shrinks the \omterm{def:g1:light-and-shadows:shadow}{shadow} back to the toy's size.
\end{solution}

\begin{solution}{exo:g1:light-and-shadows:6}
No. A \omterm{def:g1:light-and-shadows:shadow}{shadow} is not a thing made of anything --- it is a place where
\omterm{def:g1:light-and-shadows:source}{light} is missing because a blocker keeps the \omterm{def:g1:light-and-shadows:source}{light} out. With nothing
there, there is nothing to pick up.
\end{solution}

\begin{solution}{exo:g1:light-and-shadows:7}
The Sun is a far stronger source than any lamp: its \omterm{def:g1:light-and-shadows:source}{light} is so
strong it can hurt the eyes for good, not just dazzle them for a
moment.
\end{solution}

\begin{solution}{exo:g1:light-and-shadows:8}
In the late afternoon, when the Sun hangs low. The drawing should
show a short \omterm{def:g1:light-and-shadows:shadow}{shadow} at the foot of the tree at midday and a long
stretched \omterm{def:g1:light-and-shadows:shadow}{shadow} in the afternoon, both pointing away from the Sun.
\end{solution}

\begin{solution}{exo:g1:light-and-shadows:9}
Sam is right. The \omterm{def:g1:light-and-shadows:shadow}{shadow} stayed exactly opposite the Sun the whole
time. Mia turned, so the \omterm{def:g1:light-and-shadows:shadow}{shadow} changed its place \emph{compared to
her own front and back} --- but on the ground it never moved.
\end{solution}
